% TOP_PROBLEM: {"id": 2786, "title": "Kirby Problem 2.38", "category_id": 7, "category": {"id": 7, "name": "topology", "display_name": "Topology", "description": "Properties preserved under continuous deformations.", "slug": "topology", "order_index": 7, "created_at": "2026-07-31T15:26:25.671Z"}, "status": "open"}
% FIRST_POSTED: null
% ENTRY_KIND: statement_only
% Imported from: batch11.tex; UnsolvedMath ID 2786
\documentclass[11pt]{article}
\usepackage[margin=25mm]{geometry}
\usepackage{fontspec}
\setmainfont{Latin Modern Roman}
\setsansfont{Latin Modern Sans}
\usepackage{amsmath,amssymb,mathtools}
\usepackage{unicode-math}
\setmathfont{Latin Modern Math}
\usepackage{xurl,hyperref}
\setcounter{secnumdepth}{0}
\setlength{\parindent}{0pt}
\setlength{\parskip}{5pt}
\setlength{\emergencystretch}{3em}
\newcommand{\sref}[2]{\hyperlink{source:#1:#2}{[S#2]}}
\newcommand{\cataloguescope}[1]{\paragraph{Recorded target.}#1}
\begin{document}
% UnsolvedMath ID 2786; source catalogue code KP-2.38
\setcounter{section}{2785}
\section{Kirby Problem 2.38}\label{2786}
{\small\sffamily UnsolvedMath ID 2786 · Topology}
\subsection{Definitions and mathematical statement}

\paragraph{Shared notation.}$\mathbb N=\{1,2,\ldots\}$, and $\mathbb Z,\mathbb Q,\mathbb R,\mathbb C$ denote the integers, rationals, reals and complex numbers. Any other conventions are specified locally.


An oriented surface $S$ is of infinite type if it is not a compact finite-genus surface with finitely many punctures and boundary circles. Its mapping class group $\operatorname{Mod}(S)$ consists of orientation-preserving homeomorphisms modulo isotopy, fixing boundary pointwise. An end is an element of the inverse limit of the sets of components of complements of compact subsets. An end is planar if it has a genus-zero neighborhood; otherwise it is accumulated by genus. The usual curve graph has vertices essential nonperipheral simple closed curves up to isotopy, with disjointness edges. For a graph give each edge length one. Gromov hyperbolicity means that all geodesic triangles are $\delta$-thin for some finite $\delta$. An isometric action is cobounded if some orbit lies within a fixed finite distance of every point of the graph.
\paragraph{Statement recorded in the supplied source.}
Characterize exactly the infinite-type surfaces $S$ admitting an analogue of the curve graph: an associated connected infinite-diameter Gromov-hyperbolic graph on which $\operatorname{Mod}(S)$ acts isometrically and coboundedly. Construct an applicable graph for every surface in that class, and characterize the surfaces for which no such graph exists. \sref{2786}{2}

\subsection{Short English statement}
Determine which infinite-type surfaces have an infinite hyperbolic graph with a cobounded mapping-class-group action, and construct a curve-graph analogue throughout that class.
\subsection{Sources}
\begin{description}
\item[\hypertarget{source:2786:1}{S1}] ULAM.AI and the UnsolvedMath Contributors, UnsolvedMath: A Curated Collection of Open Mathematics Problems, record 2786 (KP-2.38). CC BY 4.0. \url{https://huggingface.co/datasets/ulamai/UnsolvedMath}.
\item[\hypertarget{source:2786:2}{S2}] R. İnanç Baykur, Robion C. Kirby and Daniel Ruberman, eds., \emph{K3: A New Problem List in Low-Dimensional Topology}, Mathematical Surveys and Monographs 295, American Mathematical Society (2026), Problem 2.38 and its remarks; Chapters 1 and 2 for the stated conventions. \url{https://math.berkeley.edu/sites/default/files/surv-295-ruberman-watermarked-author-pdf.pdf}.
\end{description}
\label{end:2786}
\end{document}
